\documentclass[11pt]{article}
\usepackage[a4paper,margin=0.65in]{geometry}
\usepackage[T1]{fontenc}
\usepackage{amsmath,amssymb,amsfonts}
\usepackage{graphicx}
\usepackage{pdfpages}
\usepackage[english,shorthands=off]{babel}
\usepackage{fancyhdr}
\usepackage[bookmarks=false,colorlinks=false]{hyperref}
\usepackage[protrusion=true,expansion=false]{microtype}
\emergencystretch=3em
\tolerance=600
\providecommand{\modd}{\mathrm{mod}}
\providecommand{\Disc}{\mathrm{Disc}}
\providecommand{\canont}{}
\providecommand{\asterism}{*}
\providecommand{\Hessian}{H}
\providecommand{\tome}{}
\providecommand{\page}{p.}
\pagestyle{fancy}\fancyhf{}\fancyhead[L]{Pages 163--175}\fancyhead[R]{Cayley --- Collected Papers, Vol. III}\renewcommand{\headrulewidth}{0.4pt}
\setlength{\headheight}{15pt}

\allowdisplaybreaks

\begin{document}

%===============================================================
% PDF p.163 = book p.141 -- Paper 191 (continuation)
%===============================================================
\begin{center}\small 191]\hfill\textsc{note on the expansion of the true anomaly.}\hfill 141\end{center}

\medskip

\noindent or neglecting all the terms except the first, this is
\[
= r^{r} e^{-r} \sqrt{2r} \int_{0}^{\infty} e^{-z^{2}}\, dz
\]
\[
= \sqrt{2\pi r}\; r^{r} e^{-r}.
\]

Hence multiplying by $\dfrac{1}{2^{r}}\, e^{r} e^{r} \dfrac{1}{\Gamma(r+1)}$ and observing that when $r$ is large, we have, by a well-known formula,
\[
\Gamma(r+1) = \sqrt{2\pi r}\; r^{r} e^{-r},
\]
we obtain finally the result that when $r$ is large the first term of $A_{r}$ is approximately
\[
= \left(\frac{ec}{2}\right)^{r}.
\]

I take the opportunity of mentioning the following somewhat singular theorem, which seems to belong to a more general theory: viz.\ if $u - e\sin u = m$, then we have
\[
\log(1 - e\cos u) = \frac{1}{a}\log(1 - ae\cos \phi),
\]
where
\[
\phi - \frac{1}{a}\tan\phi = m,
\]
provided that the negative powers of $a$ are rejected, and $a$ is then put equal to unity.

To show this, we have by Lagrange's theorem, observing that
\[
\frac{d}{dm} F(1 - e\cos u) = e\sin u\, F'(1 - e\cos u),
\]
\begin{align*}
F(1 - e\cos u) &= F(1 - e\cos m) + \frac{e^{2}}{1}\sin^{2} m\, F'(1 - e\cos m) \\
&\quad + \frac{e^{3}}{1\cdot 2}\, \frac{d}{dm}\sin^{3} m\, F'(1 - e\cos m) + \&\text{c.,}
\end{align*}
and the coefficient of $e^{r}$ in $F(1 - e\cos u)$ is
\[
\frac{(-)^{r}}{1\cdot 2\dots (r-1)}\left\{ \frac{1}{r}\, F_{r}\cos^{r} m + \frac{r-1}{1}\, F_{r-1}\cos^{r-2} m\sin^{2} m \right.
\]
\[
\left. - \frac{(r-1)(r-2)}{1\cdot 2}\, F_{r-2}\, \frac{d}{dm}(\cos^{r-3} m\sin^{3} m) + \&\text{c.} \right\},
\]
where $F_{r} = F^{(r)}(1)$.

\bigskip

\begin{center}\small 142\hfill\textsc{note on the expansion of the true anomaly.}\hfill [191\end{center}

\medskip

\noindent Hence in particular when $Fx = \log x$, $F_{r} = (-)^{r-1}\, 1\cdot 2\dots (r-1)$ and thence the coefficient of $e^{r}$ in $\log(1 - e\cos u)$ is
\[
- \left\{ \frac{1}{r}\cos^{r} m - \frac{1}{1\cdot 2}\cos^{r-2} m\sin^{2} m - \frac{1}{1\cdot 2}\, \frac{d}{dm}(\cos^{r-3} m\sin^{3} m) - \&\text{c.} \right\},
\]
continued as long as the exponent of $\cos$ is not negative. Now in the expansion of $\dfrac{1}{a}\log(1 - ae\cos\phi)$, where $\phi - \dfrac{1}{a}\tan\phi = m$, the coefficient of $e^{r}$ is $-\dfrac{1}{r}\, a^{r-1}\cos^{r}\phi$, and this (by Lagrange's theorem) is equal to
\[
- \frac{1}{r}\, a^{r-1}\!\left\{\cos^{r} m - \frac{1}{1\cdot 2}\, r\cos^{r-1} m\sin m\tan m - \frac{1}{1\cdot 2\cdot a^{2}}\, \frac{d}{dm}(r\cos^{r-1} m\sin m\tan m) - \&\text{c.}\right\}
\]
\[
= - \left\{ \frac{1}{r}\, a^{r-1}\cos^{r} m - \frac{1}{1}\, a^{r-3}\cos^{r-2} m\sin^{2} m - \frac{1}{1\cdot 2}\, a^{r-3}\cos^{r-3} m\sin^{3} m - \&\text{c.}\right\},
\]
where the series is continued indefinitely; but if we reject the negative powers of $a$ and then put $a$ equal to unity this is precisely equal to the former expression for the coefficient of $e^{r}$, and the formula is thus shown to be true.

\medskip

\noindent\emph{2, Stone Buildings, W.C., 17th Nov., 1857.}

\bigskip

%===============================================================
% PDF p.165 = book p.143 -- Paper 192 begins
%===============================================================
\begin{center}\small 192]\hfill\hfill 143\end{center}

\bigskip
\bigskip

\begin{center}
{\Large\bfseries 192.}\\[10pt]
{\large ON THE AREA OF THE CONIC SECTION REPRESENTED BY\\
THE GENERAL TRILINEAR EQUATION OF THE SECOND\\
DEGREE.}\\[6pt]
\small [From the \emph{Quarterly Mathematical Journal}, vol.~\textsc{ii}.\ (1858), pp.~248--253.]
\end{center}

\bigskip

\noindent [The original title was ``Direct Investigation of the Question discussed in the Foregoing Paper,'' viz.\ a Paper with the present title by N. M. Ferrers (now Dr Ferrers), pp.~247--248. The area $S$ of the conic section represented by the general equation $(A,B,C,A',B',C' \mid x,y,z)^{2} = 1$, where the coordinates are connected by the equation $x + y + z = 1$, was by considerations founded on the form of the function found to be
\[
S = \resizebox{0.86\textwidth}{!}{$\displaystyle \frac{2\pi (AA'^{2} + BB'^{2} + CC'^{2} - ABC - 2A'B'C')\, \Delta}{\left[A'^{2} - BC + B'^{2} - CA + C'^{2} - AB + 2(B'C' - AA') + 2(C'A' - BB') + 2(A'B' - CC')\right]^{\frac{3}{2}}}$}
\]
where $\Delta$ is the area of the fundamental triangle: and it was remarked that a similar method might be applied to determine the area of the conic section when it is defined by the distances of its several tangents from three given points.]

The position of a point $P$ being determined as in the foregoing paper, let $\alpha, \beta, \gamma$ denote in like manner the coordinates of a point $O$, we have
\[
\alpha + \beta + \gamma = 1,
\]
and consequently if $\xi, \eta, \zeta$ are the relative coordinates $x - \alpha,\, y - \beta,\, z - \gamma$, we have
\[
\xi + \eta + \zeta = 0.
\]

\bigskip

\begin{center}\small 144\hfill\textsc{on the area of the conic section represented by the}\hfill [192\end{center}

\medskip

The expression for the distance of the two points $O, P$ is readily obtained in terms of the relative coordinates, viz.\ calling this distance $r$, we have
\[
r^{2} = L\xi^{2} + M\eta^{2} + N\zeta^{2},
\]
where, if $l, m, n$ are the sides of the triangle $ABC$, we have
\begin{align*}
L &= \tfrac{1}{2}(m^{2} + n^{2} - l^{2}),\\
M &= \tfrac{1}{2}(n^{2} + l^{2} - m^{2}),\\
N &= \tfrac{1}{2}(l^{2} + m^{2} - n^{2});
\end{align*}
and it is to be remarked that these values give
\[
MN + LM + NL = \tfrac{1}{4}\left(2m^{2}n^{2} + 2n^{2}l^{2} + 2l^{2}m^{2} - l^{4} - m^{4} - n^{4}\right) = 4\Delta^{2},
\]
if $\Delta$ denote the area of the triangle $ABC$.

Consider now a conic
\[
(a, b, c, f, g, h \mid x, y, z)^{2},
\]
and suppose as usual that $\mathfrak{A}, \mathfrak{B}, \mathfrak{C}, \mathfrak{F}, \mathfrak{G}, \mathfrak{H}$ are the inverse coefficients and that $K$ is the discriminant, suppose also for shortness
\[
P = K\, (\mathfrak{A}, \mathfrak{B}, \mathfrak{C}, \mathfrak{F}, \mathfrak{G}, \mathfrak{H} \mid 1, 1, 1)^{2}.
\]

The coordinates of the centre being $\alpha, \beta, \gamma$, we have
\begin{align*}
\alpha &= \frac{1}{P}(\mathfrak{A}, \mathfrak{H}, \mathfrak{G} \mid 1, 1, 1),\\
\beta &= \frac{1}{P}(\mathfrak{H}, \mathfrak{B}, \mathfrak{F} \mid 1, 1, 1),\\
\gamma &= \frac{1}{P}(\mathfrak{G}, \mathfrak{F}, \mathfrak{C} \mid 1, 1, 1),
\end{align*}
and writing as before $\xi, \eta, \zeta$ for $x - \alpha,\, y - \beta,\, z - \gamma$, so that $\xi, \eta, \zeta$ are the coordinates of a point $P$ of the conic, in relation to the centre, we have $x, y, z$ respectively equal to $\xi + \alpha,\, \eta + \beta,\, \zeta + \gamma$, and the equation of the conic gives
\[
(a, \dots \mid \xi, \eta, \zeta)^{2} + 2(a, \dots \mid \alpha, \beta, \gamma)(\xi, \eta, \zeta) + (a, \dots \mid \alpha, \beta, \gamma)^{2} = 0,
\]
which may be written
\[
(a, \dots \mid \xi, \eta, \zeta)^{2}
\]
\[
+\, 2(a, \dots \mid \xi, \eta, \zeta)(\xi, \eta, \zeta)
\]
\[
+\, (a, \dots \mid \alpha, \beta, \gamma)^{2} = 0.
\]

\bigskip

\begin{center}\small 192]\hfill\textsc{general trilinear equation of the second degree.}\hfill 145\end{center}

\medskip

\noindent Now observing the equations
\begin{align*}
(a, h, g \mid \alpha, \beta, \gamma) &= \frac{K}{P},\\
(h, b, f \mid \alpha, \beta, \gamma) &= \frac{K}{P},\\
(g, f, c \mid \alpha, \beta, \gamma) &= \frac{K}{P},
\end{align*}
we have
\begin{align*}
(a, \dots \mid \alpha, \beta, \gamma)(\xi, \eta, \zeta) &= \frac{K}{P}(\xi + \eta + \zeta) = 0,\\
(a, \dots \mid \alpha, \beta, \gamma)^{2} &= \frac{K}{P}(\alpha + \beta + \gamma) = \frac{K}{P},
\end{align*}
and the equation of the conic gives therefore
\[
(a, \dots \mid \xi, \eta, \zeta)^{2} + \frac{K}{P} = 0,
\]
and we have as before
\[
\xi + \eta + \zeta = 0.
\]

To find the axes we have only to make
\[
r^{2} = L\xi^{2} + M\eta^{2} + N\zeta^{2},
\]
a maximum or minimum, $\xi, \eta, \zeta$ varying subject to the preceding two conditions; this gives
\begin{align*}
(a, h, g \mid \xi, \eta, \zeta) + \lambda L\xi + \mu &= 0,\\
(h, b, f \mid \xi, \eta, \zeta) + \lambda M\eta + \mu &= 0,\\
(g, f, c \mid \xi, \eta, \zeta) + \lambda N\zeta + \mu &= 0,
\end{align*}
and multiplying by $\xi, \eta, \zeta$, adding and reducing, we have
\[
-\frac{K}{P} + \lambda r^{2} = 0,
\]
which gives
\[
\lambda = \frac{K}{P r^{2}}.
\]

Substituting this value, and joining to the resulting three equations the equation
\[
\xi + \eta + \zeta = 0,
\]

\bigskip

\hfill c. \textsc{iii}.\hfill 19

\bigskip

\begin{center}\small 146\hfill\textsc{on the area of the conic section represented by the}\hfill [192\end{center}

\medskip

\noindent we may eliminate $\xi, \eta, \zeta, \mu$, and the result is
\[
\left| \begin{array}{cccc} a + \dfrac{KL}{P r^{2}}, & h, & g, & 1 \\[6pt] h, & b + \dfrac{KM}{P r^{2}}, & f, & 1 \\[6pt] g, & f, & c + \dfrac{KN}{P r^{2}}, & 1 \\[6pt] 1, & 1, & 1, & \end{array} \right| = 0,
\]
which may also be written
\[
(\mathfrak{A}', \mathfrak{B}', \mathfrak{C}', \mathfrak{F}', \mathfrak{G}', \mathfrak{H}' \mid 1, 1, 1)^{2} = 0,
\]
where $(\mathfrak{A}', \dots)$ are what $(\mathfrak{A}, \dots)$ become when $a, b, c$ are changed into
\[
a + \frac{KL}{P r^{2}},\quad b + \frac{KM}{P r^{2}},\quad c + \frac{KN}{P r^{2}};
\]
we in fact have
\[
\mathfrak{A}' = \mathfrak{A} + \frac{K}{P r^{2}}(bN + cM) + \frac{K^{2}}{P^{2} r^{4}}\, MN,
\]
\[
\vdots
\]
\[
\mathfrak{F}' = \mathfrak{F} - \frac{K}{P r^{2}}\, Lf,
\]
\[
\vdots
\]
and (observing the value of $P$) the result consequently is
\[
P + \frac{K}{P r^{2}}\left\{(b + c - 2f) L + (c + a - 2g) M + (a + b - 2h) N\right\} + \frac{K^{2}}{P^{2} r^{4}}(MN + NL + LM) = 0,
\]
which may also be written
\[
P^{3} r^{4} + P K r^{2}\left\{(b + c - 2f) L + (c + a - 2g) M + (a + b - 2h) N\right\} + 4\Delta^{2} K^{2} = 0.
\]

Hence if $r_{1}, r_{2}$ are the two semiaxes, we have
\[
r_{1}^{2} r_{2}^{2} = \frac{4\Delta^{2} K^{2}}{P^{3}},
\]
and the area is $\pi r_{1} r_{2}$, which is equal to
\[
\frac{2\pi K \Delta}{\sqrt{(P^{3})}},
\]
which agrees with Mr Ferrers' result.

The formula $r^{2} = L\xi^{2} + M\eta^{2} + N\zeta^{2}$ which is assumed in the preceding investigation may be proved as follows:---

\bigskip

\begin{center}\small 192]\hfill\textsc{general trilinear equation of the second degree.}\hfill 147\end{center}

\medskip

Writing $a, b, c$ (instead of $l, m, n$) for the sides of the fundamental triangle and $A, B, C$ for the angles, the equation in question is
\[
r^{2} = bc\cos A\, \xi^{2} + ca\cos B\, \eta^{2} + ab\cos C\, \zeta^{2}.
\]

Now writing $\alpha, \beta, \gamma$ for the inclinations of the line $r$ to the sides of the triangle, we have
\begin{align*}
A &= \beta - \gamma,\\
B &= \gamma - \alpha,\\
C &= \pi + \alpha - \beta.
\end{align*}

Moreover taking for a moment $\lambda, \mu, \nu$ to denote the perpendiculars from the angles on the opposite sides, we have
\begin{align*}
\lambda &= c\sin B = b\sin C,\\
\mu &= a\sin C = c\sin A,\\
\nu &= b\sin A = a\sin B,
\end{align*}
and
\[
\xi = \frac{r\sin\alpha}{\lambda},\quad \eta = \frac{r\sin\beta}{\mu},\quad \zeta = \frac{r\sin\gamma}{\nu};
\]
the values of $\xi^{2}, \eta^{2}, \zeta^{2}$ consequently are
\[
\frac{r^{2}\sin^{2}\alpha}{bc\sin B\sin C},\quad \frac{r^{2}\sin^{2}\beta}{ca\sin C\sin A},\quad \frac{r^{2}\sin^{2}\gamma}{ab\sin A\sin B},
\]
and the equation to be proved becomes
\[
1 = \frac{\cos A\sin^{2}\alpha}{\sin B\sin C} + \frac{\cos B\sin^{2}\beta}{\sin C\sin A} + \frac{\cos C\sin^{2}\gamma}{\sin A\sin B},
\]
or, what is the same thing,
\[
\sin A\sin B\sin C = \sin A\cos A\sin^{2}\alpha + \sin B\cos B\sin^{2}\beta + \sin C\cos C\sin^{2}\gamma,
\]
or again
\[
4\sin A\sin B\sin C = \sin 2A(1 - \cos 2\alpha) + \sin 2B(1 - \cos 2\beta) + \sin 2C(1 - \cos 2\gamma),
\]
or putting for $A, B, C$ their values in terms of $\alpha, \beta, \gamma$ this is
\[
-4\sin(\beta - \gamma)\sin(\gamma - \alpha)\sin(\alpha - \beta) = \sin(2\beta - 2\gamma)(1 - \cos 2\alpha)
\]
\[
+\, \sin(2\gamma - 2\alpha)(1 - \cos 2\beta) + \sin(2\alpha - 2\beta)(1 - \cos 2\gamma),
\]

\hfill 19---2

\bigskip

\begin{center}\small 148\hfill\textsc{on the area of the conic section \&c.}\hfill [192\end{center}

\medskip

\noindent which is an identical equation; it is most readily proved by writing $x, y, z$ for $\tan\alpha, \tan\beta, \tan\gamma$; the equation thus becomes
\[
\frac{-4}{(1 + x^{2})(1 + y^{2})(1 + z^{2})}(y - z)(z - x)(x - y)
\]
\[
= \Sigma \frac{1}{(1 + y^{2})(1 + z^{2})}\left[2y(1 - z^{2}) - 2z(1 - y^{2})\right]\frac{2x^{2}}{1 + x^{2}},
\]
or multiplying out
\[
-(y - z)(z - x)(x - y) = \Sigma(y - z)(1 + yz)\, x^{2} = \Sigma x^{2}(y - z) + xyz\, \Sigma x(y - z),
\]
that is
\[
-(y - z)(z - x)(x - y) = \Sigma x^{2}(y - z) = x^{2}(y - z) + y^{2}(z - x) + z^{2}(x - y),
\]
which is an identity.

\bigskip

\noindent [A different investigation of the formula $r^{2} = L\xi^{2} + M\eta^{2} + N\zeta^{2}$, by Dr Ferrers, was appended to the original Paper.]

\bigskip

%===============================================================
% PDF p.171 = book p.149 -- Paper 193 begins
%===============================================================
\begin{center}\small 193]\hfill\hfill 149\end{center}

\bigskip
\bigskip

\begin{center}
{\Large\bfseries 193.}\\[10pt]
{\large ON RODRIGUES' METHOD FOR THE ATTRACTION OF\\
ELLIPSOIDS.}\\[6pt]
\small [From the \emph{Quarterly Mathematical Journal}, vol.~\textsc{ii}.\ (1858), pp.~333--337.]
\end{center}

\bigskip

\noindent\textsc{The} following is in substance the method given in the ``M\'{e}moire sur l'attraction des Sph\'{e}ro\"{i}des,'' par M. Rodrigues, \emph{Corresp.\ sur l'\'{E}cole Polyt.}, t.~\textsc{iii}.\ pp.~361--385 (1815). It will be seen that the method is very similar to that given two years before by Gauss, see my paper ``On Gauss' Method for the Attraction of Ellipsoids,'' \emph{Journal}, t.~\textsc{i}.\ pp.~162--166 [164]: the solution in fact depends upon the geometrical theorem therein quoted, viz.\ if $M$ be any point, $P$ a point of a closed surface, $PQ$ the normal (lying outside the surface) at the point $P$, $dS$ the element of the surface at that point, and if $MQ$ denotes the angle $MPQ$ and $\overline{MP}$ the distance of the points $M$ and $P$, then, theorem, the integral
\[
\iint \frac{dS\cos MQ}{MP^{2}}
\]
has for its value
\[
0,\ 2\pi\ \text{or}\ 4\pi
\]
according as $M$ is exterior to, upon, or interior to the surface.

Suppose that $M$ is the attracted point and taking $A, B, C$ for the semiaxes of the surface of the attracting ellipsoid, or, if we please, for any semiaxes of an arbitrary ellipsoidal surface confocal with the surface of the attracting ellipsoid, let $P$ be a point on the surface of the interior similar ellipsoid whose semiaxes are $rA, rB, rC$. The coordinates of $M$ are taken to be $a, b, c$, and those of $P$ are taken to be $x, y, z$, and the value of the potential is
\[
V = \int\frac{dm}{MP},
\]
where $dm$ is the element of mass.

\bigskip

\begin{center}\small 150\hfill\textsc{on rodrigues' method for the attraction of ellipsoids.}\hfill [193\end{center}

\medskip

\noindent We may write
\begin{align*}
x &= rA\xi,\\
y &= rB\eta,\\
z &= rC\zeta,
\end{align*}
and then $\xi, \eta, \zeta$ will be the coordinates of a point $P'$ on the surface of a sphere, radius unity, corresponding in a definite manner to the point $P$ on the surface of the internal similar ellipsoid. And if $d\sigma$ be the element of the spherical surface, then we have
\[
dm = ABC\, r^{2}\, dr\, d\sigma,
\]
and therefore
\[
\frac{V}{ABC} = \int\!\!\int \frac{r^{2}\, dr\, d\sigma}{MP},
\]
where, in order to obtain the value of the potential $V$ for the ellipsoid whose semiaxes are $A, B, C$, the integrations must be extended over the spherical surface and from $r = 0$ to $r = 1$.

Suppose that $dS$ is the element of the internal similar surface at $P$, and let $p$ be the perpendicular from the centre upon the tangent plane at $P$, we have
\[
dS = \frac{r^{2} ABC}{p}\, d\sigma.
\]

Let $P_{0}$ be the point on the ellipsoid $(A, B, C)$ similarly situated to the point $P$ on the ellipsoid $(rA, rB, rC)$; the coordinates of $P_{0}$ are $A\xi, B\eta, C\zeta$; and if $p_{0}$ be the perpendicular from the centre upon the tangent plane at $P_{0}$, then $p = r p_{0}$, and the preceding equation becomes
\[
dS = \frac{r^{2} ABC}{p_{0}}\, d\sigma.
\]

Imagine now an ellipsoidal surface confocal with the surface $(A, B, C)$ and having for its semiaxes
\[
A + \delta A,\ B + \delta B,\ C + \delta C,
\]
and let $P_{0}'$ be the point on this surface which corresponds with the point $P_{0}$ on the surface $(A, B, C)$; that is, let $P_{0}'$ be the point whose coordinates are
\[
(A + \delta A)\, \xi,\ (B + \delta B)\, \eta,\ (C + \delta C)\, \zeta,
\]
and let $P'$ be in like manner the point whose coordinates are
\[
r(A + \delta A)\, \xi,\ r(B + \delta B)\, \eta,\ r(C + \delta C)\, \zeta;
\]

\bigskip

\begin{center}\small 193]\hfill\textsc{on rodrigues' method for the attraction of ellipsoids.}\hfill 151\end{center}

\medskip

\noindent the points $P', P_{0}'$ will be in like manner corresponding points on the surface $(rA, rB, rC)$ and on the confocal surface $\{r(A + \delta A), r(B + \delta B), r(C + \delta C)\}$; and if the normal distance at the point $P_{0}$ of the first two surfaces is $\delta N$, then the normal distance at the point $P$ of the second two surfaces will be $r\delta N$. The decrement of $MP$ will be equal to the normal distance $r\delta N$ of the two surfaces at the point $P$ multiplied into the cosine of the angle $MQ$, and we have, by a property of confocal surfaces,
\[
A\delta A = B\delta B = C\delta C = p_{0}\, \delta N = (\text{suppose})\ \tfrac{1}{2}\delta\theta,
\]
so we have therefore
\[
d\, MP = -\frac{\tfrac{1}{2} r\delta\theta}{p_{0}}\cos MQ.
\]

Hence from the equation
\[
\frac{V}{ABC} = \int \frac{r^{2}\, dr\, d\sigma}{MP},
\]
we deduce
\[
\delta\, \frac{V}{ABC} = \int r^{2}\, dr\, d\sigma\, \frac{\tfrac{1}{2} r\delta\theta}{p_{0}}\, \frac{\cos MQ}{MP^{2}}.
\]

But we have
\[
\frac{r^{2}\, d\sigma}{p_{0}} = \frac{dS}{ABC};
\]
and the equation thus becomes
\[
\delta\, \frac{V}{ABC} = \frac{\tfrac{1}{2}\delta\theta}{ABC}\int r\, dr\, \frac{dS\cos MQ}{MP^{2}}.
\]

It may be proper to remark here by way of recapitulation that the course of the investigation has been as follows: viz.\ that, with a view to obtaining the potential $V$ of an attracting ellipsoid, we have found the increment of $\dfrac{V}{ABC}$ in passing from the ellipsoidal surface $(A, B, C)$ to the ellipsoidal surface $(A + \delta A, B + \delta B, C + \delta C)$, each of them confocal with the surface of the attracting ellipsoid; and that for finding such increment we have had to consider the two surfaces $(rA, rB, rC)$ and $\{r(A + \delta A), r(B + \delta B), r(C + \delta C)\}$ confocal to each other and respectively similar to the first-mentioned two surfaces.

Resuming the formula just obtained, the integral with respect to $dS$ is taken over the entire surface of the internal similar ellipsoid $(rA, rB, rC)$, and if the attracted point $M$ is external to the ellipsoid $(A, B, C)$ it will be external to the interior similar ellipsoid $(rA, rB, rC)$: hence in this case the double integral vanishes for all values of $r$, or we have
\[
\delta\, \frac{V}{ABC} = 0;
\]

\bigskip

\begin{center}\small 152\hfill\textsc{on rodrigues' method for the attraction of ellipsoids.}\hfill [193\end{center}

\medskip

\noindent that is the function $V \div ABC$, which represents the ratio of the potential to the mass, is not altered in passing from the ellipsoid $(A, B, C)$ to the confocal ellipsoid
\[
(A + \delta A, B + \delta B, C + \delta C),
\]
or, what is the same thing, the potentials (and therefore the attractions) of confocal ellipsoids upon the same external point are proportional to their masses; this is in fact Maclaurin's theorem for the attraction of ellipsoids upon an external point.

But if the attracted point $M$ is interior to the ellipsoid $(A, B, C)$, then writing
\[
\frac{a^{2}}{A^{2}} + \frac{b^{2}}{B^{2}} + \frac{c^{2}}{C^{2}} = r'^{2},
\]
where $r'$ is less than unity, the double integral is $= 0$ from $r = 0$ to $r = r'$ and is $= 4\pi$ from $r = r'$ to $r = 1$, and we have
\[
\delta\, \frac{V}{ABC} = -\frac{2\pi\delta\theta}{ABC}\int_{r'}^{1} r\, dr
\]
\[
= -\frac{\pi\delta\theta}{ABC}\, (1 - r'^{2})
\]
\[
= -\frac{\pi\delta\theta}{ABC}\left(\frac{a^{2}}{A^{2}} + \frac{b^{2}}{B^{2}} + \frac{c^{2}}{C^{2}} - 1\right);
\]
that is, the right-hand side of the equation is the increment (or taken with its sign reversed so as to be positive, it is the decrement) of the function $V \div ABC$ in passing from the ellipsoid $(A, B, C)$ to the confocal ellipsoid
\[
(A + \delta A, B + \delta B, C + \delta C),
\]
where
\[
\tfrac{1}{2}\delta\theta = A\delta A = B\delta B = C\delta C.
\]

The preceding formula gives at once the potential for an interior point; in fact taking $\alpha, \beta, \gamma$ for the semiaxes of the ellipsoid and writing
\[
A^{2} = \alpha^{2} + \theta,\quad B^{2} = \beta^{2} + \theta,\quad C^{2} = \gamma^{2} + \theta,
\]
and using the ordinary symbol $d$ instead of $\delta$, we have
\[
\frac{d}{d\theta}\, \frac{V}{\sqrt{(\alpha^{2} + \theta)(\beta^{2} + \theta)(\gamma^{2} + \theta)}} = \frac{\pi}{\sqrt{(\alpha^{2} + \theta)(\beta^{2} + \theta)(\gamma^{2} + \theta)}}\left\{\frac{a^{2}}{\alpha^{2} + \theta} + \frac{b^{2}}{\beta^{2} + \theta} + \frac{c^{2}}{\gamma^{2} + \theta} - 1\right\},
\]
and integrating from $\theta = 0$ to $\theta = \infty$ we have
\[
-\frac{V}{\alpha\beta\gamma} = \pi\int_{0}^{\infty} \frac{d\theta}{\sqrt{(\alpha^{2} + \theta)(\beta^{2} + \theta)(\gamma^{2} + \theta)}}\left\{\frac{a^{2}}{\alpha^{2} + \theta} + \frac{b^{2}}{\beta^{2} + \theta} + \frac{c^{2}}{\gamma^{2} + \theta} - 1\right\},
\]

\bigskip

\begin{center}\small 193]\hfill\textsc{on rodrigues' method for the attraction of ellipsoids.}\hfill 153\end{center}

\medskip

\noindent where $V$ is now the potential for the ellipsoid whose semiaxes are $\alpha, \beta, \gamma$; and we have therefore
\[
V = -\pi\alpha\beta\gamma\int_{0}^{\infty} \frac{d\theta}{\sqrt{(\alpha^{2} + \theta)(\beta^{2} + \theta)(\gamma^{2} + \theta)}}\left\{\frac{a^{2}}{\alpha^{2} + \theta} + \frac{b^{2}}{\beta^{2} + \theta} + \frac{c^{2}}{\gamma^{2} + \theta} - 1\right\}.
\]

To find the potential for an external point it is only necessary to remark that by the theorem above demonstrated $V \div \alpha\beta\gamma$ is equal to the corresponding function for the confocal ellipsoid through the attracted point, that is for the ellipsoid whose semiaxes are $\sqrt{(\alpha^{2} + \theta_{1})}, \sqrt{(\beta^{2} + \theta_{1})}, \sqrt{(\gamma^{2} + \theta_{1})}$, where $\theta_{1}$ is a positive quantity such that
\[
\frac{a^{2}}{\alpha^{2} + \theta_{1}} + \frac{b^{2}}{\beta^{2} + \theta_{1}} + \frac{c^{2}}{\gamma^{2} + \theta_{1}} = 1;
\]
hence in the value of $V \div \alpha\beta\gamma$ we have only to write the above values in the place of $\alpha, \beta, \gamma$; and if we then write $\theta - \theta_{1}$ in the place of $\theta$ the limits will be $\infty, \theta_{1}$, and the expression for the potential is
\[
V = -\pi\alpha\beta\gamma\int_{\theta_{1}}^{\infty} \frac{d\theta}{\sqrt{(\alpha^{2} + \theta)(\beta^{2} + \theta)(\gamma^{2} + \theta)}}\left\{\frac{a^{2}}{\alpha^{2} + \theta} + \frac{b^{2}}{\beta^{2} + \theta} + \frac{c^{2}}{\gamma^{2} + \theta} - 1\right\};
\]
this completes the investigation.

\bigskip

\hfill c. \textsc{iii}.\hfill 20

\end{document}
